\chapter[Koszul--Petersson and the Prime Gibbs Model][Prime Gibbs model and pairings]{Koszul--Petersson and the Prime Gibbs Model}
\label{v4-arith:kms-gns-full}

A single prime admits an explicit Gibbs state for every inverse
temperature \(\beta>0\), including \(\beta=1\). Its GNS space is a
Hilbert--Schmidt space, its modular conjugation is adjunction, and
its modular operator has a maximal weighted-matrix domain. These
objects can be computed without identifying an arithmetic Koszul
pairing with a Petersson pairing. The latter comparison requires
additional maps between the actual carriers.

\section{The Prime State and the Arithmetic Comparison}
\label{v4-arith:kms-gns:statement-table}
\label{v4-arith:kms-gns:residuals}
\label{v4-arith:kms-gns:consequences}
\label{v4-arith:kms-gns:P3-consequence}
\label{v4-arith:kms-gns:table}

The occupation model of Chapter~\ref{v4-arith:prim-zeta-char} has
Hamiltonian \(j\log p\) at a prime. We first construct its GNS
representation and an actual unitary onto the prime-generated
cyclic subspace of any specified global KMS representation.
Neither this subspace comparison nor uniqueness of a polar
decomposition supplies an automorphic realization.

\begin{hypothesis}[arithmetic placement and pairing comparison]
\label{v4-arith:kms-gns:hyp:arithmetic-pairing}
Fix a primitive arithmetic carrier, its coefficient field
\(\mathbb C\), its completion and a dense test domain. Specify:
\begin{enumerate}
\item an injective placement \(\iota\) of the selected domain into a
fixed-central-character automorphic Hilbert space, with its Haar
and Tamagawa normalizations;
\item at each place, comparison maps for the independently defined
Koszul and automorphic carriers, intertwining the stated algebra
or Hecke actions and transporting the bilinear Koszul form and
antilinear involution to the stated local Hermitian form;
\item preservation of the chosen reference vectors at almost all
places, factorization on elementary tensors, and continuous or
explicitly regularized extension of the resulting pairing identity;
\item on every included residual sector, its specified placement and
normalization, and orthogonality of the sector decomposition for
both transported forms.
\end{enumerate}
Any modular comparison must additionally intertwine the represented
algebra, state or cyclic vector, and the closed modular operators
on their actual domains. Ramified sectors require their own local
maps; they are not consequences of the prime shift computation.
\end{hypothesis}

\begin{question}[constructing the arithmetic pairing maps]
\label{v4-arith:kms-gns:question:arithmetic-pairing}
Construct these maps from the arithmetic bar--cobar carrier and
prove their compatibility with the independently defined Petersson
pairing. Determine their precise discrete and continuous domains.
The conditional transfers below retain this question; the prime
Gibbs calculation does not complete it.
\end{question}

\section{The One-Prime Gibbs State}
\label{v4-arith:kms-gns:cycle1}
\subsection{State and Physical Time}
\label{v4-arith:kms-gns:cycle1:state}

Fix a prime \(p\) and \(\beta>0\). Put
\(K=\ell^2(\mathbb N_0)\), \(Ve_j=e_{j+1}\), and
\(He_j=j\log p\,e_j\), with
\(\operatorname{Dom}H=\{x:\sum_jj^2|x_j|^2<\infty\}\).
Let \(\mathcal T=C^*(V)\subset B(K)\),
\(q=p^{-\beta}\), \(r_j=(1-q)q^j\), and
\(\rho=(1-q)q^{H/\log p}\). Physical time is
\begin{equation}
\label{v4-arith:kms-gns:eq:sigma-time}
\alpha_t(a)=e^{itH}ae^{-itH},\qquad \alpha_t(V)=p^{it}V.
\end{equation}
The positive injective trace-class operator \(\rho\) has trace one.
Its faithful normal state on \(B(K)\), restricted to \(\mathcal T\), is
\begin{equation}
\label{v4-arith:kms-gns:eq:phi-1}
\varphi_{p,\beta}(a)=\operatorname{Tr}(\rho a),\qquad
\varphi_{p,\beta}(V^mV^{*n})=\delta_{m,n}q^m.
\end{equation}
In particular \(\varphi_{p,\beta}(V^*V)=1\) and
\(\varphi_{p,\beta}(VV^*)=q\). The order of the isometry and its
adjoint is essential.

\begin{lemma}[prime KMS state]
\label{v4-arith:kms-gns:lem:KMS-property}
The state \(\varphi_{p,\beta}\) is the unique \(\beta\)-KMS state
on \((\mathcal T,\alpha)\), with convention
\begin{equation}
\label{v4-arith:kms-gns:eq:KMS}
\varphi_{p,\beta}(ab)=\varphi_{p,\beta}(b\alpha_{i\beta}(a))
\end{equation}
for analytic elements.
\end{lemma}
\begin{proof}
The analytic monomials \(V^mV^{*n}\) span a dense star algebra.
On them \(\rho a=\alpha_{i\beta}(a)\rho\), so trace cyclicity
proves the KMS identity. Conversely, applying that identity with
\(b=1\) forces the expectation of a monomial of nonzero degree
to vanish. Taking \(a=V^m,b=V^{*m}\) gives expectation \(q^m\)
for \(V^mV^{*m}\), since \(V^{*m}V^m=1\). Thus every KMS state
has the values in \eqref{v4-arith:kms-gns:eq:phi-1}; density gives
uniqueness.
\end{proof}

\subsection{The Hilbert--Schmidt GNS Carrier}
\label{v4-arith:kms-gns:cycle1:GNS}

Use inner products linear in the first variable.
\begin{lemma}[GNS realization of the prime state]
\label{v4-arith:kms-gns:lem:GNS-space}
The GNS representation is
\begin{equation}
\label{v4-arith:kms-gns:eq:GNS-iso}
\mathcal H_{p,\beta}=\operatorname{HS}(K),\quad
\pi(a)X=aX,\quad \Omega=\rho^{1/2},
\end{equation}
with
\begin{equation}
\label{v4-arith:kms-gns:eq:GNS-inner-product}
\langle X,Y\rangle=\operatorname{Tr}(XY^*),\qquad
[a]\longmapsto a\rho^{1/2}.
\end{equation}
The vector \(\Omega\) is cyclic and separating for left \(B(K)\),
which is \(\pi(\mathcal T)''\).
\end{lemma}
\begin{proof}
The displayed map is isometric because its inner product is
\(\operatorname{Tr}(\rho b^*a)=\varphi_{p,\beta}(b^*a)\).
The matrix units \(E_{ij}=V^i(1-VV^*)V^{*j}\) belong to
\(\mathcal T\), and \(E_{ij}\rho^{1/2}=\sqrt{r_j}E_{ij}\).
Their images span a dense subspace of \(\operatorname{HS}(K)\).
Finite matrix corners strongly approximate bounded operators,
so the weak closure is left \(B(K)\). If \(a\Omega=0\), every
column of \(a\) vanishes because every \(r_j>0\); hence \(a=0\).
\end{proof}

\subsection{Closed Modular Operators}
\label{v4-arith:kms-gns:cycle1:modular}

Define the densely defined antilinear operator
\begin{equation}
\label{v4-arith:kms-gns:eq:S-def}
S_0(a\Omega)=a^*\Omega,\qquad a\in\mathcal T.
\end{equation}
\begin{lemma}[modular operator and maximal domains]
\label{v4-arith:kms-gns:lem:Delta}
For \(X=\sum x_{ij}E_{ij}\), set \(JX=X^*\) and
\begin{equation}
\label{v4-arith:kms-gns:eq:Delta-basis}
\Delta E_{ij}=q^{i-j}E_{ij},\qquad
\operatorname{Dom}(\Delta^a)=
\left\{X:\sum_{i,j}q^{2a(i-j)}|x_{ij}|^2<\infty\right\},\quad a\in\mathbb R.
\end{equation}
Then \(\overline{S_0}=J\Delta^{1/2}\), on the displayed domain for
\(a=1/2\), and
\begin{equation}
\label{v4-arith:kms-gns:eq:Delta-formula}
\Delta^{it}X=e^{-i\beta tH}Xe^{i\beta tH}.
\end{equation}
Thus the modular action on left multiplication is physical time
\(\alpha_{-\beta t}\). Both \(\Delta\) and \(\Delta^{-1}\) are
unbounded. The generator \(\log\Delta\) has eigenvalue
\((i-j)\log q\) and domain defined by square summability with
weight \(|(i-j)\log q|^2\).
\end{lemma}
\begin{proof}
On matrix units, \(S_0E_{ij}=q^{(i-j)/2}E_{ji}\), antilinearly.
Finite matrix sums are graph-dense for the maximal weighted
operator, by truncation of the convergent double sum. For any
bounded \(a\), the matrix \(a\Omega\) lies in its domain and the
formula gives \(a^*\Omega\); its weighted norm is
\(\|a^*\Omega\|_{\mathrm{HS}}<\infty\). Therefore the closure is
exactly \(J\Delta^{1/2}\). Diagonal spectral calculus gives all
powers, domains and the unitary group. The eigenvalues \(q^{i-j}\)
range through \(q^{\mathbb Z}\), proving both unboundedness claims.
\end{proof}

\begin{proposition}[modular conjugation and the prime shift]
\label{v4-arith:kms-gns:prop:J-equals-sigma-BC}
For \(k\geq0\),
\begin{align}
\label{v4-arith:kms-gns:eq:J-basis}
J(V^k\Omega)&=p^{\beta k/2}V^{*k}\Omega,\\
\label{v4-arith:kms-gns:eq:J-on-algebra}
JL_aJ&=R_{a^*},\qquad a\in B(K),
\end{align}
where \(L\) and \(R\) denote left and right multiplication.
The continuous bilinear form \(B(X,Z)=\operatorname{Tr}(XZ)\)
on \(\operatorname{HS}(K)^2\) satisfies
\(B(X,JY)=\langle X,Y\rangle\).
\end{proposition}
\begin{proof}
Expand \(V^k\Omega=\sum_j\sqrt{r_j}E_{j+k,j}\) and use
\(\sqrt{r_{j+k}}=q^{k/2}\sqrt{r_j}\). Adjunction gives the first
identity and \(J L_aJ(X)=Xa^*\) gives the second.
Products of two Hilbert--Schmidt operators are trace class with
\(|\operatorname{Tr}(XZ)|\leq\|X\|_2\|Z\|_2\), proving continuity
and the asserted pairing identity. In particular the half-power
has positive exponent: \(\|V^k\Omega\|=1\), whereas
\(\|V^{*k}\Omega\|=q^{k/2}\).
\end{proof}

\begin{theorem}[comparison with a prime cyclic subspace]
\label{v4-arith:kms-gns:thm:KMS1-GNS-BC-identification}
Let \((\mathcal A,\alpha,\psi_\beta)\) be a specified dynamical
\(C^*\)-algebra with a \(\beta\)-KMS state and an isometry
\(\mu_p\) satisfying \(\alpha_t(\mu_p)=p^{it}\mu_p\).
Let \(\iota_p:\mathcal T\to\mathcal A\) send \(V\) to \(\mu_p\).
There is a canonical unitary
\[
U_p:\operatorname{HS}(K)\longrightarrow
\overline{\pi_{\psi_\beta}(C^*(\mu_p))\Omega_{\psi_\beta}},\qquad
U_p(a\rho^{1/2})=\pi_{\psi_\beta}(\iota_p(a))\Omega_{\psi_\beta}.
\]
It intertwines the prime algebra actions and cyclic vectors, and
transports the closed modular operators of the restricted cyclic
representation. This holds at \(\beta=1\) as well.
\end{theorem}
\begin{proof}
The pulled-back state \(\psi_\beta\circ\iota_p\) is a KMS state
on \(\mathcal T\), hence equals \(\varphi_{p,\beta}\) by the
unique-state calculation. The displayed map preserves inner
products on \(\mathcal T\rho^{1/2}\) and has dense image in the
specified cyclic subspace. It therefore extends to a unitary.
It intertwines multiplication and the operation
\(a\Omega\mapsto a^*\Omega\) on its core, hence its closure and
polar decomposition.
\end{proof}

\begin{remark}[scope of modular uniqueness]
\label{v4-arith:kms-gns:rem:uniqueness}
This unitary targets the prime cyclic subspace, not the whole
\(\psi_\beta\)-GNS space. It does not identify the phase operators,
an automorphic representation or an arithmetic Koszul form.
Uniqueness of modular conjugation on a fixed cyclic representation
supplies no additional identification between those carriers.
\end{remark}

\subsection{Finite Prime Sets and the Critical Global Boundary}
\label{v4-arith:kms-gns:cycle1:adelic}
\begin{corollary}[finite-prime Gibbs and GNS products]
\label{v4-arith:kms-gns:cor:adelic-KMS-GNS}
For a finite prime set \(F\), the density
\(\rho_F=\bigotimes_{p\in F}\rho_p\) has trace one for every
\(\beta>0\). Its GNS space and modular operators are the finite
tensor products of the preceding constructions. On the full
occupation space \(\ell^2(\mathbb N^{\times})\), the heat trace is
\(\zeta(\beta)\) for \(\beta>1\), and diverges at \(\beta=1\).
Thus the full occupation representation has no normalized
trace-class Gibbs density at the critical temperature.
\end{corollary}
\begin{proof}
Finite tensor products multiply traces and tensor the explicit
matrix formulas. Unique prime factorization identifies occupations
with integers and gives \(\sum_{n\geq1}n^{-\beta}\). At
\(\beta=1\) this is the divergent harmonic series.
\end{proof}
The divergence does not preclude a critical Bost--Connes KMS state
in another representation. For any such specified state the
prime cyclic comparison above still applies. Neither argument
identifies its entire GNS space or factor type with the finite-prime
Gibbs model.

\section{Local Pairing Models and Arithmetic Transport}
\label{v4-arith:kms-gns:cycle2}
\subsection{Ramified Principal-Series Comparison}
\label{v4-arith:kms-gns:cycle2:ramified}

\begin{lemma}[a normalized local Hermitian carrier]
\label{v4-arith:kms-gns:lem:GJ-ramified}
Fix Haar measure \(d^\times a\) on \(\mathbb Q_p^\times\).
On \(E_p=C_c^\infty(\mathbb Q_p^\times,\mathbb C)\), put
\begin{equation}
\label{v4-arith:kms-gns:eq:GJ-ramified}
 I_p(v,w)=\int_{\mathbb Q_p^\times}v(a)\overline{w(a)}
 |a|^{-1}\,d^\times a.
\end{equation}
This is positive definite and extends to the corresponding
\(L^2\)-completion. The bilinear form
\(b_p(v,z)=\int v(a)z(a)|a|^{-1}\,d^\times a\) satisfies
\(b_p(v,\overline w)=I_p(v,w)\).
\end{lemma}
\begin{proof}
Compact support in \(\mathbb Q_p^\times\) gives finite integrals.
The measure is positive, and a nonzero locally constant function
is nonzero on an open set of positive measure. Cauchy--Schwarz
extends the forms to \(L^2\), and complex conjugation gives the
last identity.
\end{proof}

The Whittaker realization of a given representation, its scalar
normalization and its identification with an arithmetic Koszul
carrier are further data. In particular an epsilon factor from a
local functional equation is not automatically a positive
normalization of a Hermitian form.

\begin{hypothesis}[local Koszul comparison]
\label{v4-arith:kms-gns:lem:Koszul-ramified}
For the chosen ramified arithmetic domain \(V_{p,\pi}\), specify a
linear map \(c_{p,\pi}\) into the local Hilbert carrier, a compatible
map on its Koszul dual, a positive constant \(c_\pi\), and the
antilinear map \(\sigma_p\). Require that the actual Koszul pairing
satisfy
\begin{equation}
\label{v4-arith:kms-gns:eq:Koszul-ramified}
 \langle v,\sigma_p(w)\rangle^{\mathrm{Koszul}}_{p,\pi}
 =c_\pi b_p(c_{p,\pi}v,\overline{c_{p,\pi}w}).
\end{equation}
The maps must preserve the specified local actions, domains and
completion. This is an arithmetic comparison hypothesis.
\end{hypothesis}

\begin{hypothesis}[local automorphic normalization]
\label{v4-arith:kms-gns:lem:Pet-ramified}
For the same carrier, maps, measure and positive constant, require
\begin{equation}
\label{v4-arith:kms-gns:eq:Pet-ramified}
 \langle\iota_p v,\iota_p w\rangle^{\mathrm{Pet}}_{p,\pi}
 =c_\pi I_p(c_{p,\pi}v,c_{p,\pi}w).
\end{equation}
Neither Whittaker uniqueness nor a scalar local functional equation
alone identifies the map \(c_{p,\pi}\) with the Koszul comparison.
\end{hypothesis}

\begin{proposition}[conditional local Koszul--Petersson comparison]
\label{v4-arith:kms-gns:prop:local-match-ramified}
Under the two local comparison hypotheses above,
\(\langle\iota_pv,\iota_pw\rangle^{\mathrm{Pet}}_{p,\pi}
=\langle v,\sigma_p(w)\rangle^{\mathrm{Koszul}}_{p,\pi}\)
on the specified domain.
\end{proposition}
\begin{proof}
Substitute \(b_p(x,\overline y)=I_p(x,y)\) into the two displayed
identities. The common constant and measure are part of the
hypotheses, not a conclusion of matching the names of the pairings.
\end{proof}

\subsection{Supercuspidal Comparison}
\label{v4-arith:kms-gns:cycle2:supercuspidal}
\begin{lemma}[conditional supercuspidal transport]
\label{v4-arith:kms-gns:lem:supercuspidal}
The same conclusion holds on a supercuspidal domain for which the
local Koszul and automorphic comparison hypotheses have been
constructed with a common Hilbert carrier and normalization.
\end{lemma}
\begin{proof}
The proof of the preceding proposition uses only the common forms
and maps, so applies on that specified domain. It supplies no new
comparison map for the supercuspidal representation.
\end{proof}

\subsection{Temperedness and Assembly}
\label{v4-arith:kms-gns:cycle2:ramanujan}
\begin{proposition}[conditional extension under temperedness]
\label{v4-arith:kms-gns:prop:non-tempered-residue}
Assume Hypothesis~\ref{v4-arith:kms-gns:hyp:arithmetic-pairing} on a
tempered cuspidal domain. Any further cuspidal domain contained in
it inherits its pairing comparison. Applying this to Maass forms
requires the appropriate Ramanujan temperedness assumption as well
as the arithmetic comparison maps.
\end{proposition}
\begin{proof}
Restrict the specified maps and pairing equality to the subdomain.
Temperedness is a domain condition and does not construct the maps.
\end{proof}

\begin{remark}[local exponent bounds]
\label{v4-arith:kms-gns:rem:partial-KS}
The Kim--Sarnak exponent bound \(\theta\leq7/64\) is a local
bound. It is neither an arithmetic pairing map nor a summable
prime error on the critical line. The exponent \(-25/64\) gives
\(\sum_p p^{-25/64}=\infty\); it cannot be used as an absolutely
convergent adelic error majorant.
\end{remark}

\begin{theorem}[conditional cuspidal pairing assembly]
\label{v4-arith:kms-gns:thm:ramified-cuspidal-tempered-extension}
Assume Hypothesis~\ref{v4-arith:kms-gns:hyp:arithmetic-pairing},
including the local comparisons at every ramified place and the
specified adelic extension. On the resulting tempered cuspidal
domain,
\begin{equation}
\label{v4-arith:kms-gns:eq:ramified-cuspidal-match}
 \langle\iota v,\iota w\rangle^{\mathrm{Pet}}
 =\langle v,\sigma^{\mathrm{BC}}(w)\rangle^{\mathrm{Koszul}}.
\end{equation}
\end{theorem}
\begin{proof}
On elementary restricted tensors the forms factor, and the local
identities give equality of every finite product with the specified
reference-vector normalization. The assumed continuous or
regularized extension carries this equality to its stated domain.
Finiteness of the ramified set alone does not justify the infinite
product normalization or its extension.
\end{proof}

\begin{remark}[scope of the domain]
\label{v4-arith:kms-gns:rem:w3-tightening}
The normalized local Hilbert models and the finite-prime Gibbs
models are explicit. Enlarging an arithmetic domain to ramified,
residual or Maass sectors requires its actual comparison maps.
The theorem above is their conditional assembly, not an
unconditional enlargement of the arithmetic pairing theorem.
\end{remark}

\section{Langlands Decomposition of \(L^{2}([GL_{2}])\) and Placement of \(\mathcal{V}^{\mathrm{prim}}\)}
\label{v4-arith:kms-gns:cycle3}

\subsection{The Spectral Decomposition}
\label{v4-arith:kms-gns:cycle3:spectral}

The Langlands decomposition of \(L^{2}([GL_{2}])\), following
Langlands 1966 and Moeglin--Waldspurger 1995, is:

\begin{theorem}[spectral decomposition of \(L^{2}(\lbrack GL_{2}\rbrack)\)]
\label{v4-arith:kms-gns:thm:spectral-decomposition}
There is an orthogonal decomposition of
\(G(\mathbb{A})\)-modules:
\begin{equation}
\label{v4-arith:kms-gns:eq:L2-decomposition}
L^{2}([GL_{2}])\;=\;L^{2}_{\mathrm{cusp}}([GL_{2}])\;\oplus\;L^{2}_{\mathrm{res}}([GL_{2}])\;\oplus\;L^{2}_{\mathrm{cont}}([GL_{2}]),
\end{equation}
where:
\begin{itemize}
\item \(L^{2}_{\mathrm{cusp}}([GL_{2}])\) is the cuspidal spectrum,
a direct sum of irreducible unitary representations
\(\pi=\bigotimes'_{v}\pi_{v}\) with \(\pi_{p}\) generic tempered at
almost all primes;
\item \(L^{2}_{\mathrm{res}}([GL_{2}])\) is the residual spectrum
(in the \(GL_{2}\) case generated by one-dimensional characters
pulled up from the determinant), consisting of \(GL_{2}(\mathbb{A})\)-reducible
automorphic representations from residues of Eisenstein series at
specific points;
\item \(L^{2}_{\mathrm{cont}}([GL_{2}])\) is the continuous spectrum,
realised as the direct integral
\(\int^{\oplus}_{\mathrm{Re}(s)=1/2}\mathrm{Ind}_{B}^{G}(\chi_{s})\,|\xi(2s)|^{-2}ds\)
for \(\chi_{s}=|\cdot|^{s}\otimes|\cdot|^{-s}\).
\end{itemize}
\end{theorem}

\begin{proof}
Langlands 1966 \emph{On the Functional Equations Satisfied by
Eisenstein Series} \S6--7 for the construction; Moeglin--Waldspurger
1995 \emph{Spectral Decomposition and Eisenstein Series for \(GL_{n}\)}
for the published \(GL_{n}\) exposition; for \(GL_{2}\) specifically see
Gelbart 1975 \emph{Automorphic Forms on Adele Groups} Ch.~II.
\end{proof}

\subsection{Placement of \(\mathcal{V}^{\mathrm{prim}}\)}
\label{v4-arith:kms-gns:cycle3:placement}

\begin{hypothesis}[discrete arithmetic placement]
\label{v4-arith:kms-gns:prop:Vprim-placement}
Choose the primitive arithmetic vector carrier, its dense domain
and its completion. The discrete placement required in
Hypothesis~\ref{v4-arith:kms-gns:hyp:arithmetic-pairing} is an
injective map
\begin{equation}
\label{v4-arith:kms-gns:eq:placement}
 \iota:\mathcal V^{\mathrm{prim}}_{\mathrm{disc}}
 \longrightarrow L^2_{\mathrm{cusp}}\oplus L^2_{\mathrm{res}},
\end{equation}
with fixed central character and specified Haar normalization.
Its chosen residual component and vacuum normalization are part of
the datum. The Langlands decomposition of the target does not
construct this map. A continuous-spectrum comparison requires its
own direct-integral or distributional carrier and is not a map of
a distribution into \(L^2\) without further hypotheses.
\end{hypothesis}

\section{Scalar Continuum Models and the Pairing Problem}
\label{v4-arith:kms-gns:cycle4}
\subsection{The Scalar Coefficient Domain}
\label{v4-arith:kms-gns:cycle4:maass-selberg}

Use the Riemann completed function \(\xi\) with
\(\xi(1-z)=\xi(z)\) and \(\xi(\overline z)=\overline{\xi(z)}\).
Let \(U\) be the points \(s\) on \(\operatorname{Re}s=1/2\)
where \(\xi(2s)\) and \(\xi(2s-1)\) are finite and nonzero.

\begin{definition}[scalar scattering coefficient model]
\label{v4-arith:kms-gns:lem:maass-selberg}
For \(s\in U\), set \(d(s)=|\xi(2s)|^{-2}>0\) and
\(M(s)=\xi(2s-1)/\xi(2s)\). Define
\begin{equation}
\label{v4-arith:kms-gns:eq:maass-selberg}
 P(s)=(1+|M(s)|^2)d(s).
\end{equation}
This is a scalar model. Identifying it with a coefficient of an
actual truncated Petersson pairing requires that pairing, its
truncation parameter and its measure normalization.
\end{definition}

\subsection{The Proposed Koszul Coefficient}
\label{v4-arith:kms-gns:cycle4:koszul-continuation}
\begin{definition}[scalar Koszul coefficient model]
\label{v4-arith:kms-gns:lem:Koszul-continued}
On the same domain define
\begin{equation}
\label{v4-arith:kms-gns:eq:Koszul-continued}
 K(s)=M(s)^{-1}d(s)
 =\frac{\xi(2s)}{\xi(2s-1)}|\xi(2s)|^{-2}.
\end{equation}
This scalar is not asserted to be the spectral measure of an
arithmetic Koszul pairing. Such a comparison requires the
continuous carrier, coefficient map and regularization.
\end{definition}

\begin{theorem}[strict discrepancy of the two scalar models]
\label{v4-arith:kms-gns:thm:eisenstein-discrepancy}
For every \(s\in U\), one has \(|M(s)|=1\), \(P(s)=2d(s)\), and
\begin{equation}
\label{v4-arith:kms-gns:eq:continuum-discrepancy}
 P(s)-K(s)=d(s)(2-M(s)^{-1}).
\end{equation}
In particular
\begin{equation}
\label{v4-arith:kms-gns:eq:continuum-match-condition}
 d(s)\leq |P(s)-K(s)|\leq3d(s),
\end{equation}
so the models are unequal everywhere on \(U\).
\end{theorem}
\begin{proof}
Write \(s=1/2+it\). Functional equation and conjugation give
\(\xi(2s)=\xi(1+2it)=\xi(-2it)=\overline{\xi(2it)}\).
The numerator and denominator of \(M\) therefore have equal
nonzero modulus. Substitution gives the difference, and the
triangle and reverse triangle inequalities give the two bounds.
\end{proof}

\begin{remark}[phase and scale]
\label{v4-arith:kms-gns:rem:match-is-FE}
Multiplication by the phase \(M\) gives \(MK=d=P/2\), not \(P\).
Multiplication by \(2M\) matches these scalar models. Neither
scalar identity constructs a positive spectral measure, an
isometry, or an arithmetic bar--cobar comparison.
\end{remark}

\begin{corollary}[scope of the scalar obstruction]
\label{v4-arith:kms-gns:cor:eisenstein-final}
A proposed continuous pairing comparison whose coefficients are
exactly \(P\) and \(K\) on \(U\), with the same measure and no
additional transport, cannot identify them. A comparison using
another normalization must construct that normalization and its
maps; the scalar discrepancy does not exclude such a construction.
\end{corollary}
\begin{proof}
Their difference has modulus at least \(d>0\). The last assertion
only separates this fixed-normalization obstruction from a
comparison with different data.
\end{proof}

\begin{remark}[continuous pairings require more than one coefficient]
\label{v4-arith:kms-gns:rem:continuum-honest-domain}
The scalar functions above specify no bivariate pairing kernel.
In particular they do not determine an off-diagonal truncated
Eisenstein pairing or its diagonal limit. On the critical line,
\(s_1+\overline{s_2}-1=i(t_1-t_2)\), which vanishes exactly when
\(s_1=s_2\). Any formula dividing by this quantity must treat that
diagonal with its actual truncation and limit. The classical
Maass--Selberg and Arthur identities concern the automorphic
side; they supply no integrated equality with an arithmetic
Koszul pairing without the missing comparison maps.
\end{remark}

\section{Conditional Arithmetic Pairing Consequences}
\label{v4-arith:kms-gns:P3}

\begin{theorem}[conditional P3 on the specified cuspidal domain]
\label{v4-arith:kms-gns:thm:P3-tempered-cuspidal}
Assume Hypothesis~\ref{v4-arith:kms-gns:hyp:arithmetic-pairing} and
the discrete placement of
Hypothesis~\ref{v4-arith:kms-gns:prop:Vprim-placement} on the chosen
tempered cuspidal carrier. Then the pullback of the Petersson form
is positive definite and agrees there with the transported Koszul
form composed with the specified antilinear involution.
\end{theorem}
\begin{proof}
Injectivity of \(\iota\) makes the pullback of a positive-definite
form positive definite. The local pairing comparisons and their
specified adelic extension give equality of forms by
Theorem~\ref{v4-arith:kms-gns:thm:ramified-cuspidal-tempered-extension}.
The prime GNS theorem alone does not supply either hypothesis.
\end{proof}

\begin{remark}[domain of arithmetic transport]
\label{v4-arith:kms-gns:rem:P3-domain-enlargement}
The prime Gibbs and Hilbert--Schmidt calculations are unconditional
on their stated carriers. Extension of the arithmetic pairing to
ramified, holomorphic, Maass or residual sectors requires its
placement and pairing maps on each sector. Ramanujan controls a
temperedness condition; it is not the missing isometry. Continuous
spectrum additionally requires its measure comparison.
\end{remark}

\subsection{The Sufficient RH Implication}
\label{v4-arith:kms-gns:P3:rh}
\begin{corollary}[conditional use in a sufficient RH implication]
\label{v4-arith:kms-gns:cor:rh-minimal-updated}
A sufficient RH implication that assumes P3 may use the preceding
conditional theorem only on its specified domain and with all of
its arithmetic comparison hypotheses. Identifying the relevant
Li or Weil trace with that domain is a further spectral comparison.
The other hypotheses of the sufficient implication remain in force.
\end{corollary}
\begin{proof}
Substitute the conditional pairing conclusion for the P3 premise
only after its hypotheses and the shared trace domain are supplied.
This logical substitution constructs neither set of hypotheses.
\end{proof}

\section{Prime Models and Remaining Maps}
\label{v4-arith:kms-gns:residual-list}
\label{v4-arith:kms-gns:verification}

The positive constructions are the normalized prime Gibbs state,
its Hilbert--Schmidt GNS realization, the closed modular operators,
the finite-prime product, and the unitary onto a specified global
state's prime cyclic subspace. Their domains and state
normalizations are explicit. The local integral form
\eqref{v4-arith:kms-gns:eq:GJ-ramified} is another independently
defined positive Hermitian model.

The arithmetic construction still requires the maps in
Hypothesis~\ref{v4-arith:kms-gns:hyp:arithmetic-pairing}. The
residual-spectrum vacuum, ramified Whittaker models and continuous
spectral measures must be identified on their actual carriers.
A scalar partition function, the uniqueness of modular conjugation,
and a spectral decomposition of the target do not provide those
maps.

\section{The Arithmetic Comparison Problem}
\label{v4-arith:kms-gns:summary}

Question~\ref{v4-arith:kms-gns:question:arithmetic-pairing} retains
the intended Koszul--Petersson comparison. The finite and prime
models specify normalization tests and actual Hilbert-space maps;
the arithmetic placement and pairing compatibility remain to be
constructed. At the critical global temperature, the prime cyclic
comparison continues to hold for a specified KMS state even though
the full occupation Gibbs trace diverges.
